% Part: first-order-logic
% Chapter: completeness
% Section: identity

\documentclass[../../../include/open-logic-section]{subfiles}

\begin{document}

\olfileid{fol}{com}{ide}
\olsection{एकरूपता}

\begin{explain}
मागील विभागात दिलेली पद-प्रतिमानाची रचना ही~$\eq$ नसलेल्या
संच~$\Gamma$ साठी प्रथम-कोटी तर्कशास्त्राची संपूर्णता प्रस्थापित करण्यास
पुरेशी आहे. पद-प्रतिमान हे~$\eq$ नसलेले प्रत्येक $!A \in \Gamma^*$
पूर्त करते (आणि म्हणून सर्व $!A \in \Gamma$ पूर्त करते). पण~$\eq$
उपस्थित असेल, तर ही रचना चालत नाही. कारण मग~$\Gamma^*$ मध्ये
!!a{sentence}~$\eq[t][t']$ असू शकते, पण पद-प्रतिमानात कोणत्याही पदाचे
मूल्य ते पद स्वतःच असते. म्हणून $t$ आणि $t'$ ही भिन्न पदे असतील, तर
पद-प्रतिमानातील त्यांची मूल्ये---म्हणजे अनुक्रमे $t$ आणि $t'$---भिन्न
असतात; त्यामुळे $\eq[t][t']$ असत्य असते. मात्र ``भागाकार घेणे'' म्हणून
ओळखली जाणारी रचना वापरून हे दुरुस्त करता येते.
\end{explain}

\begin{defn}
  $\Gamma^*$ हा~$\Lang L$ मधील वाक्यांचा सुसंगत आणि !!{complete} संच
  असू द्या. $\Lang L$ मधील बंद पदांच्या संचावर आपण~$\approx$ हा संबंध
  पुढीलप्रमाणे परिभाषित करतो:
  \[
  t \approx t' \text{\quad तेव्हा आणि केवळ तेव्हाच \quad} \eq[t][t'] \in \Gamma^*
  \]
\end{defn}

\begin{prop}
\ollabel{prop:approx-equiv}
$\approx$ या संबंधाला पुढील गुणधर्म आहेत:
\begin{enumerate}
\item $\approx$ परावर्ती आहे.
\item $\approx$ सममित आहे.
\item $\approx$ संक्रमक आहे.
\item $t \approx t'$, $f$ हे !!a{function}, आणि $t_1$, \dots,
  $t_{i-1}$, $t_{i+1}$, \dots, $t_n$ ही बंद पदे असतील, तर
\[
\Atom{f}{t_1,\dots, t_{i-1}, t, t_{i+1}, \dots, t_n} \approx
\Atom{f}{t_1,\dots, t_{i-1}, t', t_{i+1}, \dots, t_n}.
\]
\item $t \approx t'$, $R$ हे !!a{predicate}, आणि $t_1$, \dots,
  $t_{i-1}$, $t_{i+1}$, \dots, $t_n$ ही बंद पदे असतील, तर
\begin{multline*}
\Atom{R}{t_1,\dots, t_{i-1}, t, t_{i+1}, \dots, t_n} \in \Gamma^* \text{ तेव्हा आणि केवळ तेव्हाच } \\
\Atom{R}{t_1,\dots, t_{i-1}, t', t_{i+1}, \dots, t_n} \in \Gamma^*.
\end{multline*}
\end{enumerate}
\end{prop}

\begin{proof}
$\Gamma^*$ सुसंगत आणि !!{complete} असल्यामुळे $\eq[t][t'] \in
\Gamma^*$ तेव्हा आणि केवळ तेव्हाच, जेव्हा $\Gamma^* \Proves
\eq[t][t']$. म्हणून पुढील विधाने दाखवणे पुरेसे आहे:
\begin{enumerate}
\item प्रत्येक बंद पद~$t$ साठी $\Gamma^* \Proves \eq[t][t]$.
\item $\Gamma^* \Proves \eq[t][t']$ असेल, तर $\Gamma^* \Proves \eq[t'][t]$.
\item $\Gamma^* \Proves \eq[t][t']$ आणि $\Gamma^* \Proves
  \eq[t'][t'']$ असेल, तर $\Gamma^* \Proves \eq[t][t'']$.
\item $\Gamma^* \Proves \eq[t][t']$ असेल, तर
\[
\Gamma^* \Proves
\eq[\Atom{f}{t_1,\dots,t_{i-1},t,t_{i+1},\dots,t_n}][\Atom{f}{t_1,\dots,t_{i-1},t',t_{i+1},\dots,t_n}]
\]
प्रत्येक $n$-स्थानी !!{function}~$f$ आणि बंद पदे $t_1$, \dots,
$t_{i-1}$, $t_{i+1}$, \dots,~$t_n$ यांच्यासाठी.
\item $\Gamma^* \Proves \eq[t][t']$ आणि
$\Gamma^* \Proves
\Atom{R}{t_1,\dots,t_{i-1},t,t_{i+1},\dots,t_n}$ असेल, तर
$\Gamma^* \Proves \Atom{R}{t_1,\dots,t_{i-1},t',t_{i+1},\dots,t_n}$
प्रत्येक $n$-स्थानी !!{predicate}~$R$ आणि बंद पदे $t_1$, \dots,
$t_{i-1}$, $t_{i+1}$, \dots,~$t_n$ यांच्यासाठी.
\end{enumerate}
%READERNOTE{OLCOM-004: गोठवलेल्या इंग्रजी स्रोताच्या फलन-पदातील t_{i+1} नंतर सलग दोन स्वल्पविराम आहेत. मराठीत एकच स्वल्पविराम ठेवला असून गोठवलेले इंग्रजी बाइट्स बदललेले नाहीत.}
\end{proof}

\begin{prob}
\olref[fol][com][ide]{prop:approx-equiv} ची सिद्धता पूर्ण करा.
\end{prob}

\begin{defn}
$\Gamma^*$ हा भाषा~$\Lang L$ मधील सुसंगत आणि !!{complete} संच आहे,
$t$ हे बंद पद आहे, आणि~$\approx$ हा मागील व्याख्येतील संबंध आहे असे
समजा. मग:
\[
\equivrep{t}{\approx} = \Setabs{t'}{t'\in \Trm[L], t \approx t'}
\]
आणि $\equivclass{\Trm[L]}{\approx} = \Setabs{\equivrep{t}{\approx}}{t \in \Trm[L]}$.
\end{defn}

\begin{defn}
\ollabel{defn:term-model-factor}
$\Struct M = \Struct M(\Gamma^*)$ हे \olref[mod]{defn:termmodel} मधील
$\Gamma^*$ चे पद-प्रतिमान असू द्या. मग
$\Struct{\equivclass{M}{\approx}}$ ही पुढील !!{structure} आहे:
\begin{enumerate}
\item $\Domain{\equivclass{M}{\approx}} = \equivclass{\Trm[L]}{\approx}$.
\item $\Assign{c}{\equivclass{M}{\approx}} = \equivrep{c}{\approx}$
\item $\Assign{f}{\equivclass{M}{\approx}}(\equivrep{t_1}{\approx}, \dots,
  \equivrep{t_n}{\approx}) = \equivrep{\Atom{f}{t_1,\dots, t_n}}{\approx}$
\item $\tuple{\equivrep{t_1}{\approx}, \dots, \equivrep{t_n}{\approx}} \in
  \Assign{R}{\equivclass{M}{\approx}}$ तेव्हा आणि केवळ तेव्हाच, जेव्हा
  $\Sat{M}{\Atom{R}{t_1,\dots, t_n}}$; म्हणजे तेव्हा आणि केवळ तेव्हाच,
  जेव्हा $\Atom{R}{t_1,\dots, t_n} \in \Gamma^*$.
\end{enumerate}
\end{defn}

\begin{explain}
लक्षात घ्या की~$\equivclass{\Trm[L]}{\approx}$ च्या घटकांना
$\equivrep{t}{\approx}$ असे संबोधून, म्हणजेच
\emph{प्रतिनिधी}~$t \in \equivrep{t}{\approx}$ यांच्या साहाय्याने,
आपण त्यांच्यासाठी $\Assign{f}{\equivclass{M}{\approx}}$ आणि
$\Assign{R}{\equivclass{M}{\approx}}$ परिभाषित केली आहेत. या व्याख्या
प्रतिनिधींच्या निवडीवर अवलंबून नाहीत याची आपल्याला खात्री केली पाहिजे.
म्हणजे, तेच तुल्यतावर्ग निश्चित करणारी दुसरी~$t'$ निवडली
($\equivrep{t}{\approx} = \equivrep{t'}{\approx}$), तरी व्याख्यांतून तेच
फल मिळाले पाहिजे. उदाहरणार्थ, $R$ हे एकस्थानी !!{predicate} असेल, तर
व्याख्येच्या शेवटच्या बाबीनुसार $\equivrep{t}{\approx} \in
\Assign{R}{\equivclass{M}{\approx}}$ तेव्हा आणि केवळ तेव्हाच, जेव्हा
$\Sat{M}{\Atom{R}{t}}$. आता $t \approx t'$ असलेल्या दुसऱ्या
पद~$t'$ साठी $\Sat/{M}{\Atom{R}{t'}}$ असेल, तर व्याख्येनुसार
$\equivrep{t'}{\approx} \notin \Assign{R}{\equivclass{M}{\approx}}$ असावे
लागेल. $t \approx t'$ असेल, तर $\equivrep{t}{\approx} =
\equivrep{t'}{\approx}$; पण $\equivrep{t}{\approx} \in
\Assign{R}{\equivclass{M}{\approx}}$ आणि $\equivrep{t}{\approx}
\notin \Assign{R}{\equivclass{M}{\approx}}$ ही दोन्ही विधाने लागू होऊ
शकत नाहीत. मात्र \olref{prop:approx-equiv} मुळे असे होऊ शकत नाही, याची
हमी मिळते.
%READERNOTE{OLCOM-005: गोठवलेल्या इंग्रजी स्रोताच्या प्रतिनिधी-स्वातंत्र्य उदाहरणात दुसऱ्या प्रतिनिधी t' विषयी बोलताना न-पूर्ति सूत्रात पुन्हा t आले आहे. मराठीत t' पुनर्स्थापित केले असून गोठवलेले इंग्रजी बाइट्स बदललेले नाहीत.}
\end{explain}

\begin{prop}
$\Struct{\equivclass{M}{\approx}}$ सुव्याख्यायित आहे; म्हणजे $t_1$,
  \dots, $t_n$, $t_1'$, \dots, $t_n'$ ही बंद पदे असतील आणि
  $t_i \approx t_i'$ असेल, तर
\begin{enumerate}
\item $\equivrep{\Atom{f}{t_1,\dots, t_n}}{\approx} =
    \equivrep{\Atom{f}{t_1',\dots, t_n'}}{\approx}$, म्हणजे
  \[
  \Atom{f}{t_1,\dots, t_n} \approx \Atom{f}{t_1',\dots, t_n'}
  \]
  आणि
\item $\Sat{M}{\Atom{R}{t_1,\dots, t_n}}$ तेव्हा आणि केवळ तेव्हाच,
  जेव्हा $\Sat{M}{\Atom{R}{t_1',\dots, t_n'}}$; म्हणजे
  \[
    \Atom{R}{t_1,\dots, t_n} \in \Gamma^* \text{ तेव्हा आणि केवळ तेव्हाच }
    \Atom{R}{t_1',\dots, t_n'} \in \Gamma^*.
  \]
\end{enumerate}
\end{prop}

\begin{proof}
\olref{prop:approx-equiv} वरून~$n$ वरील विगमनाने मिळते.
\end{proof}

पद-प्रतिमानाच्या बाबतीत होते त्याप्रमाणेच, सत्यता पूर्वप्रमेय सिद्ध
करण्यापूर्वी आपल्याला पुढील पूर्वप्रमेयाची गरज आहे.

\begin{lem}
\ollabel{lem:val-in-termmodel-factored} $\Struct M = \Struct M
 (\Gamma^*)$ असू द्या; मग $\Value{t}{\equivclass{M}{\approx}} = \equivrep{t}
 {\approx}$.
\end{lem}
\begin{proof}
ही सिद्धता \olref[mod]{lem:val-in-termmodel} च्या सिद्धतेसारखीच आहे.
\end{proof}

\begin{prob}
\olref[fol][com][ide]{lem:val-in-termmodel-factored} ची सिद्धता पूर्ण करा.
\end{prob}

\begin{lem}
\ollabel{lem:truth}
सर्व वाक्ये~$!A$ यांच्यासाठी $\Sat{\equivclass{M}{\approx}}{!A}$ तेव्हा आणि
केवळ तेव्हाच, जेव्हा $!A \in \Gamma^*$.
\end{lem}

\begin{proof}
\olref[mod]{lem:truth} च्या सिद्धतेप्रमाणेच~$!A$ वरील विगमनाने सिद्धता
मिळते. $!A \ident \eq[t][t']$ असण्याच्या प्रसंगाकडेच फक्त अधिक लक्ष
देण्याची गरज आहे.
\begin{align*}
\Sat{\equivclass{M}{\approx}}{\eq[t][t']} & \text{ तेव्हा आणि केवळ तेव्हाच } \equivrep{t}{\approx} = \equivrep{t'}{\approx}
\text{ ($\Struct{\equivclass{M}{\approx}}$ च्या व्याख्येनुसार)}\\
& \text{ तेव्हा आणि केवळ तेव्हाच } t \approx t' \text{ ($\equivrep{t}{\approx}$ च्या व्याख्येनुसार)}\\
& \text{ तेव्हा आणि केवळ तेव्हाच } \eq[t][t'] \in \Gamma^* \text{ ($\approx$ च्या व्याख्येनुसार).}
\end{align*}
\end{proof}

\begin{digress}
लक्षात घ्या की $\Struct{M(\Gamma^*)}$ हे नेहमी !!{enumerable} आणि अनंत
असले, तरी $\Struct{\equivclass{M}{\approx}}$ सांत असू शकते; कारण
$\equivrep{t}{\approx}$ असे फक्त सांत अनेक वर्ग असू शकतात. हे अपेक्षितच
आहे, कारण~$\Gamma$ मध्ये अशी !!{sentence}s असू शकतात की ती सत्य असणारी
कोणतीही !!{structure} सांत असणे आवश्यक असते. उदाहरणार्थ,
$\lforall[x][\lforall[y][x = y]]$ हे सुसंगत !!{sentence} आहे, पण ज्याच्या
!!a{domain} मध्ये नेमका एकच !!{element} आहे अशाच !!{structure}s मध्ये
त्याची पूर्ति होते.
\end{digress}

\end{document}
